% Supplementary project detail. This fragment inherits the preamble and
% formatting of Project_Report.tex and is included with \input.

\section{Supplementary Project Record}
This section gathers additional implementation details that are useful
when reading, reproducing, or demonstrating the project. It follows the
same report class, page dimensions, typography, line spacing, headers,
and numbering as the main report because it is included from the main
LaTeX document. The chapter documents the implemented prototype and
identifies proposed checks separately from completed work.

\subsection{Two Distinct Project Workflows}
The project has an offline model-development workflow and a separate
interactive scoring workflow. They share feature definitions, but they
have different inputs, objectives, and outputs. Keeping them distinct
is necessary to interpret the reported test metrics correctly.

\subsubsection{Offline development workflow}
The notebook reads the synthetic activity files and answer file,
constructs employee-day features, attaches known-incident-window
labels, and creates employee-grouped train, validation, and test
partitions. Candidate estimators are fitted on training data. Validation
Average Precision is used to compare ranking performance, and the
candidate review capacities are compared on validation data. The held-
out test partition is then used for the reported final estimates.

After the evaluation stage, the notebook retrains the artifact models
using the available labelled employee-day data. For the Random Forest,
the selected synthetic augmentation rule is applied during this final
fit when validation selected that variant. This full-data model bundle
is intended for interactive demonstration. The test metrics describe
the earlier held-out evaluation model and must not be presented as a
fresh evaluation of the final bundle, since that bundle has seen the
records used in the test partition during its final fit.

\subsubsection{Interactive scoring workflow}
The local dashboard accepts activity files without the answer file. It
rebuilds employee-day features, aligns the resulting columns to the
feature order saved in the model bundle, and scores every row with all
four fitted estimators. It converts each score series to a percentile
rank within the selected folder, computes a consensus value, aggregates
daily ranks to employee rankings, and writes two CSV summaries. Since
this inference workflow does not receive answer labels, it cannot
calculate precision, recall, or incident-window correctness for the
selected folder.

The distinction is summarized in Table~\ref{tab:workflow-boundaries}.

\begin{table}[H]
\centering
\caption{Offline evaluation and interactive scoring boundaries}
\label{tab:workflow-boundaries}
\small
\begin{tabularx}{\textwidth}{|p{0.22\textwidth}|X|X|}
\hline
Aspect & Offline development and evaluation & Interactive dashboard scoring\\
\hline
Activity input & CERT r4.2 logs & User-selected folder with five activity logs\\
\hline
Answer file & Used to construct known-window labels & Not required and not uploaded\\
\hline
Purpose & Compare models and estimate ranking metrics on held-out groups & Generate a relative review queue for a selected folder\\
\hline
Model fit & Train, validate, evaluate, then create a final demonstration bundle & Load saved estimators; no fitting occurs during a folder analysis\\
\hline
Evaluation target & Employee-day known-window label & No correctness target is available at inference time\\
\hline
Main output & Saved comparison metrics and test predictions & Employee ranking and employee-day score CSV files\\
\hline
Interpretation & Research estimate for a fixed split and label rule & Within-folder prioritization; not a measured probability\\
\hline
\end{tabularx}
\end{table}

\subsection{Detailed Input Contract}
The current shared scoring implementation expects five filenames and
uses normalized column names. The fields in Table~\ref{tab:input-contract}
describe the implemented feature processor. Optional columns affect
only the feature derived from that column; the basic event count still
uses each retained row. A missing required date or user field stops
processing for that source.

\begin{longtable}{|p{0.14\textwidth}|p{0.20\textwidth}|p{0.22\textwidth}|p{0.34\textwidth}|}
\caption{Input columns and derived values in the dashboard processor}\label{tab:input-contract}\\
\hline
Source & Required columns & Optional columns & Derived employee-day values\\
\hline
\endfirsthead
\hline
Source & Required columns & Optional columns & Derived employee-day values\\
\hline
\endhead
Logon & \texttt{date}, \texttt{user} & None used by the feature processor & Event count and after-hours event count\\
\hline
Device & \texttt{date}, \texttt{user} & \texttt{activity} & Event count and after-hours count; connection count when activity is ``connect''\\
\hline
File & \texttt{date}, \texttt{user} & None used by the feature processor & Event count and after-hours event count\\
\hline
Email & \texttt{date}, \texttt{user} & \texttt{attachments} & Event count and after-hours count; numeric attachment sum when available\\
\hline
HTTP & \texttt{date}, \texttt{user} & None used by the feature processor & Event count and after-hours event count\\
\hline
Optional employee lookup & A recognized user identifier column and name column & Lookup filename may be one of the supported variants & Display name attached to the employee report; no model feature is created\\
\hline
\end{longtable}

The timestamp field is parsed with invalid values coerced to missing.
Rows with a missing or blank user, an invalid date, or a duplicate within
the current chunk do not contribute to the aggregate. User identifiers
are trimmed as strings. Date values are converted to a day key, while
the hour is used only to create the fixed after-hours indicator. These
rules mean that the aggregate is not a lossless representation of the
input: the original event order and the full timestamp are not retained
as model inputs.

\subsubsection{Feature alignment}
The saved model bundle carries an ordered feature list. Before scoring,
the processor checks each expected feature and inserts a zero column
when that feature is absent from the selected aggregate. It then selects
the ordered columns and coerces values to numeric form, filling
non-numeric values with zero. This makes inference tolerant of optional
source fields, but it can also hide a mismatch between the training
schema and an incomplete source export. A stronger future implementation
could distinguish an intentionally absent optional field from a failed
or incomplete export and report both conditions explicitly.

\subsection{Review-Queue Calculation}
Let \(s_{u,d,m}\) denote the raw score for user \(u\), day \(d\), and
model \(m\). Each model's scores are ranked within the current analysis
and expressed as percentile values \(r_{u,d,m}\). The daily consensus
is
\[
c_{u,d}=\frac{1}{4}\sum_{m=1}^{4}r_{u,d,m}.
\]
The employee ranking retains the maximum value across that user's days,
\[
C_u=\max_d c_{u,d},
\]
and also retains each model's maximum daily percentile. Model agreement
counts how many of those employee-level model ranks fall in that
model's top decile. In implementation, the display status is assigned
when the consensus reaches the top-five-percent employee cutoff or
when at least two model ranks meet the top-decile condition. The report
uses ``SUSPICIOUS'' for the stored status value, while the interface
uses review-oriented wording such as ``Prioritize''; neither string is
a determination of misconduct.

\subsubsection{Signal text and its limitations}
The generated reason string is based on the employee's highest-
consensus day. The code checks a limited set of counts against the
95th percentile of the corresponding feature across all employee-days
in the current folder. A signal is included only when the feature is
present, the value is positive, and the feature percentile cutoff is
positive. At most three signal phrases are retained; if no phrase meets
the rule, a generic high-ranking activity message is used.

This is a descriptive summary rather than a model explanation. The
signal text does not show feature contributions, causal relationships,
or the sequence of raw events. The feature quantiles are recalculated
for each analysis, so a phrase may appear or disappear when the
composition of the selected folder changes. Analysts should use it as
a navigation aid to the aggregate day, then consult authorized source
records and context through a separate process.

\subsection{Illustrative Review Walkthrough}
This walkthrough describes how an analyst could navigate a result in
the demonstration. It is a procedural example and does not describe an
actual employee or a saved investigation outcome.

\begin{enumerate}
  \item The analyst starts the local API and React client and confirms
        that the health indicator reports the model bundle as available.
  \item The analyst selects a folder containing the five named source
        files. The browser reports any missing required filenames
        before enabling analysis.
  \item During submission, the interface reports upload progress and
        then backend processing milestones. The remaining-time text is
        treated as an estimate because the processing stages vary in
        duration.
  \item After a successful response, the analyst reviews the overall
        number of users, flagged count, and active-day count. These are
        descriptive totals for the selected folder, not quality
        measures for the model.
  \item The analyst inspects the highest-ranked employee cards and
        considers the four-model agreement count alongside the consensus
        rank. Disagreement can still be informative because the four
        models use different scoring principles.
  \item The analyst opens an employee profile and checks the four
        within-folder percentile ranks, the generated signal text, and
        the highest-ranked days. No raw event content is shown in this
        profile.
  \item If a lead requires further attention, the analyst checks the
        original authorized source records and relevant schedule or
        operational context, then records the outcome in an approved
        system outside this prototype.
\end{enumerate}

At each step, the user should remember that rankings depend on the
selected folder. If the folder omits dates, users, or one source file,
the resulting ranks describe only the records that were successfully
processed. A high rank does not establish intent, and a low rank does
not establish that no incident occurred.

\subsection{Proposed Acceptance and Quality Plan}
The project report records notebook and saved-artifact checks, but it
does not claim a completed end-to-end user acceptance test. The cases
below form a concrete plan for a later verification run. Every case is
marked as proposed so it is not confused with completed verification.

\begin{longtable}{|p{0.08\textwidth}|p{0.23\textwidth}|p{0.40\textwidth}|p{0.19\textwidth}|}
\caption{Supplementary acceptance and quality plan}\label{tab:supplementary-qa}\\
\hline
ID & Area & Test procedure and expected result & Status\\
\hline
\endfirsthead
\hline
ID & Area & Test procedure and expected result & Status\\
\hline
\endhead
QA01 & Required-file validation & Omit one required source and verify the API names the missing file and does not score a partial set & Proposed\\
\hline
QA02 & Header normalization & Vary header capitalization and surrounding spaces; verify normalized required headers are recognized & Proposed\\
\hline
QA03 & Invalid timestamp & Include an invalid date and verify the row does not affect daily event counts & Proposed\\
\hline
QA04 & Blank identity & Include an empty and whitespace-only user; verify both rows are excluded & Proposed\\
\hline
QA05 & Chunk boundary & Duplicate a row at the end of one chunk and start of the next; document that current chunk-local deduplication may retain both & Proposed\\
\hline
QA06 & Optional device field & Remove the activity column and verify that device event counts remain available while connection counts are absent and aligned appropriately & Proposed\\
\hline
QA07 & Optional attachment field & Remove or corrupt the attachment column and verify email event counts still process without a fabricated attachment total & Proposed\\
\hline
QA08 & Empty source & Provide a valid header and no data rows; verify the overall analysis handles an empty source alongside nonempty sources & Proposed\\
\hline
QA09 & Empty selection & Provide five empty files and verify a clear no-usable-records response & Proposed\\
\hline
QA10 & Bundle mismatch & Remove one model or the feature list from a controlled test bundle and verify loading stops with a useful diagnostic & Proposed\\
\hline
QA11 & Rank ties & Supply identical feature rows and inspect stable tie handling and resulting percentile values & Proposed\\
\hline
QA12 & Employee detail & Open a valid employee, then an unknown identifier; verify detail and not-found behaviors respectively & Proposed\\
\hline
QA13 & Output schema & Confirm employee-day and employee-level CSV columns, row counts, and user/day keys after one successful run & Proposed\\
\hline
QA14 & Retention and permissions & Inspect temporary staging cleanup and generated-output permissions after success and after a forced error & Proposed\\
\hline
QA15 & Accessibility & Complete keyboard-only navigation and screen-reader review of folder selection, progress, errors, queue, and profile & Proposed\\
\hline
\end{longtable}

The quality plan includes expected behavior that can be observed or
recalculated, rather than relying on a subjective impression that the
interface appears to work. For each completed run, a future test record
should include the input fixture version, software versions, expected
and observed result, evidence location, and reviewer. Sensitive real
activity should not be used as a test fixture unless it has been
explicitly approved and protected.

\subsection{Artifact and Run Record}
For a reproducible demonstration, a run record can be maintained
alongside the generated files. Table~\ref{tab:run-record} provides a
blank schema for such a record. No values are supplied here because
the report does not contain a documented end-to-end dashboard analysis
run.

\begin{table}[H]
\centering
\caption{Suggested run-record fields for a future demonstration}
\label{tab:run-record}
\small
\begin{tabularx}{\textwidth}{|p{0.29\textwidth}|X|}
\hline
Field & Purpose\\
\hline
Run identifier & Distinguishes one analysis from another without using a person identifier\\
\hline
Input snapshot reference & Identifies approved file set using controlled location and checksums\\
\hline
Operator and authorization reference & Records who performed the run and its approved purpose in an authorized record system\\
\hline
Model bundle version & Connects the run to the estimator and ordered feature list used\\
\hline
Processing start/end time & Supports operational timing and resource measurement\\
\hline
Quality summary & Records input rows, excluded rows, available fields, and any processing warnings\\
\hline
Output references & Identifies generated employee and day summaries under appropriate access control\\
\hline
Retention and deletion date & Records how long the run artifacts may be retained and when they are removed\\
\hline
Review outcome reference & Links to an approved external review record without storing a sensitive narrative in the model output\\
\hline
\end{tabularx}
\end{table}

This record schema is a proposal, not a feature of the present API.
Adding run identifiers or audit metadata would require coordinated
changes to the service, UI, storage permissions, and retention process.
It should be implemented only as part of a defined data-governance
design.

\subsection{Known Technical Boundaries}
The current demonstration has several boundaries that should remain
visible when its outputs are presented:
\begin{itemize}
  \item The answer-file window is a proxy label and may not identify
        which event on a marked day relates to the scenario.
  \item The employee-grouped split tests unseen synthetic users, but it
        does not establish time-forward generalization.
  \item Synthetic positive augmentation is selected on one validation
        split and is not evidence that the generated rows match real
        activity patterns.
  \item Employee-day classification metrics do not measure the
        employee-level dashboard rule or analyst outcomes.
  \item Percentile scores and queue thresholds are recalculated within
        each selected dataset and cannot be compared as absolute risk
        probabilities across runs.
  \item The application does not authenticate analysts, retain case
        decisions, or enforce authorization, access, and deletion
        policies.
  \item The captured UI images show the landing page and scoring guide;
        they do not demonstrate a completed end-to-end analysis screen.
\end{itemize}

These constraints define the appropriate scope of the project record:
an educational, local prototype for feature engineering, comparative
model evaluation, and review-queue interaction. The contents of this
report do not establish suitability for employment decisions or
real-world monitoring.
